\section{A normalized and invertible smoothness coordinate}\label{sec:smoothness}
For a fixed numeric $\mu$, define the conditional effective degrees of freedom $\nu_d(\lambda)=\tr H_\lambda$. The $d$ polynomial-nullspace eigenvalues of $H_\lambda$ equal one, and the other $r=n-d$ eigenvalues equal $(1+\lambda\delta_j)^{-1}$. Therefore
\begin{equation}\label{eq:edf}
\nu_d(\lambda)=d+\sum_{j=1}^r\frac{1}{1+\lambda\delta_j}.
\end{equation}
The classical trace reduction divided by $n$ is
\begin{equation}\label{eq:classicalS}
S^{\rm trace}_d(\lambda)=\frac{n-\nu_d(\lambda)}{n},
\quad 0\le S^{\rm trace}_d\le1-\frac dn.
\end{equation}
The normalized coordinate used here divides by the number of \emph{penalized} directions:
\begin{equation}\label{eq:S}
\boxed{\;
S_d(\lambda)=\frac{n-\nu_d(\lambda)}{n-d}
=\frac{1}{r}\sum_{j=1}^r\frac{\lambda\delta_j}{1+\lambda\delta_j}.
\;}
\end{equation}
Thus $S_d=(n/r)S_d^{\rm trace}$, an interpretable rescaling of a published index \citep{guerrero2008,cortes2017}, not a different family of fitted trends.

\begin{theorem}[Smoothness bijection]\label{thm:bijection}
For every integer $1\le d<n$, $S_d$ is continuous and strictly increasing on the nonnegative finite real line, with $S_d(0)=0$ and $\lim_{\lambda\to\infty}S_d(\lambda)=1$. It extends to a bijection between $[0,\infty]$ and $[0,1]$. Each $s\in(0,1)$ has a unique inverse $\lambda=\Lambda_{n,d}(s)$ independent of $y$ and any fixed numeric $\mu$.
\end{theorem}
\begin{proof}
All $\delta_j$ are strictly positive by Lemma~\ref{lem:kernel}, so
\begin{equation}\label{eq:Sprime}
S_d'(\lambda)=\frac{1}{r}\sum_{j=1}^r
\frac{\delta_j}{(1+\lambda\delta_j)^2}>0.
\end{equation}
Continuity and the stated endpoint values follow term by term. The inverse exists by the intermediate value theorem and strict monotonicity. Formula \eqref{eq:S} contains only $(n,d,\lambda)$ and the spectrum of $D_d$, so it is independent of the observed values and fixed reference mean.
\end{proof}

The same requested $s$ has a common interpretation across different orders, but generally corresponds to \emph{different} penalties $\Lambda_{n,d}(s)$. Direct optimization in raw $\lambda$ or $\log\lambda$ is possible. The point of $s$ is a bounded, physically interpretable search parameter, not a proof of optimization supremacy.

\begin{proposition}[Change of selection coordinate]\label{prop:coordinate}
Let $C(\lambda)$ be any fixed-$\mu$ criterion on an admissible finite interval and set $\widetilde C(s)=C(\Lambda_{n,d}(s))$. The two objectives attain the same global minimizers after the monotone change of coordinates, provided the domains and endpoint conventions match. For $s\in(0,1)$,
\begin{equation}\label{eq:chain}
\frac{d\widetilde C}{ds}
=\frac{dC/d\log\lambda}{dS_d/d\log\lambda},
\quad
\frac{dS_d}{d\log\lambda}
=\frac1r\sum_{j=1}^r\frac{\lambda\delta_j}
{(1+\lambda\delta_j)^2}.
\end{equation}
\end{proposition}
\begin{proof}
The first statement follows from the bijection. The second is the chain rule and $d\lambda/d\log\lambda=\lambda$.
\end{proof}
The denominator in \eqref{eq:chain} tends to zero at both endpoints; hence boundedness alone does not resolve numerical conditioning or spurious local minima of the criterion. Global minimizer equivalence does not imply that bounded and logarithmic numerical optimizers follow identical paths.

\paragraph{Exact endpoint conventions.}
The point $s=0$ means $\lambda=0$ and the interpolation fit $y$; the point $s=1$ means the constrained fit \eqref{eq:inftyfit}, not a floating-point stand-in for an extremely large finite penalty. Some criteria are undefined at these endpoints even though the fitted trends exist. A robust implementation must treat such limits analytically rather than blindly calling a root finder at $s=0$ or $s=1$.

\paragraph{Caching the inverse map.}
For fixed $n,d$, the full set $\{\Lambda_{n,d}(s_j)\}_{j=1}^M$ can be mapped once and reused across multiple fixed-$\mu$ values, independent series, and criteria. In contrast, an independently constructed uniform grid in $\log\lambda$ generally selects different candidate trends and is not an identical-workload control.
